% Generated by generate.py: do not edit.
\chapter{Part III: Toward Research}
\section{Tome 10: Chance II}
\subsection*{1.\ Independence and Zero--One Laws}
\begin{results}
\item[1.1] Lemma ($\pi$-systems suffice)
\item[1.2] Proposition (grouping)
\item[1.4] Theorem (Kolmogorov's zero--one law, 1933)
\item[1.7] Theorem (Hewitt--Savage, 1955)
\item[1.8] Theorem (the fate of a random walk)
\item[1.10] Lemma (Ottaviani's inequality)
\item[1.11] Theorem (L\'evy, 1937)
\item[1.12] Proposition (Kolmogorov's $L^2$ criterion)
\item[1.13] Lemma (Paley--Zygmund)
\item[1.14] Lemma (moments of a Rademacher sum)
\item[1.15] Theorem (Rademacher, 1922; Khintchine--Kolmogorov, 1925)
\end{results}
\subsection*{2.\ Laws of Large Numbers}
\begin{results}
\item[2.1] Theorem (weak law for triangular arrays)
\item[2.2] Theorem (Feller, 1945)
\item[2.3] Theorem (Etemadi, 1981)
\item[2.4] Lemma (Kolmogorov's other inequality)
\item[2.5] Theorem (three-series theorem, Kolmogorov)
\item[2.6] Lemma (Kronecker)
\item[2.7] Theorem (strong law of renewal)
\end{results}
\subsection*{3.\ Weak Convergence and Donsker's Theorem}
\begin{results}
\item[3.1] Lemma (uniqueness)
\item[3.2] Lemma (truncation inequality)
\item[3.3] Theorem (L\'evy's continuity theorem)
\item[3.4] Corollary (Cram\'er--Wold device)
\item[3.6] Proposition (continuous mapping)
\item[3.7] Lemma (converging together)
\item[3.8] Theorem (strong Markov property)
\item[3.10] Proposition (exit from an interval)
\item[3.11] Theorem (Skorokhod's embedding, 1961)
\item[3.12] Corollary (embedding a random walk)
\item[3.13] Theorem (Donsker's invariance principle, 1951)
\item[3.14] Corollary (the maximum)
\item[3.16] Lemma (Chung--Feller)
\item[3.17] Theorem (L\'evy's arcsine law, 1939; Erd\H{o}s--Kac, 1947)
\item[3.18] Theorem (law of the iterated logarithm for Brownian motion)
\item[3.19] Theorem (Hartman--Wintner, 1941)
\end{results}
\subsection*{4.\ The Central Limit Theorem, Sharpened}
\begin{results}
\item[4.1] Theorem (Lindeberg--Feller)
\item[4.2] Corollary (Lyapunov's condition)
\item[4.4] Lemma (P\'olya's density)
\item[4.5] Lemma (Esseen's smoothing inequality)
\item[4.7] Theorem (Berry--Esseen)
\item[4.9] Theorem (multivariate central limit theorem)
\item[4.12] Theorem (local limit theorem; Gnedenko, 1948)
\item[4.14] Theorem (recurrence on $\Z^d$)
\end{results}
\subsection*{5.\ Infinitely Divisible and Stable Laws}
\begin{results}
\item[5.1] Lemma (continuous logarithm)
\item[5.3] Theorem (L\'evy--Khintchine)
\item[5.5] Theorem (convergence to infinitely divisible laws)
\item[5.7] Theorem (symmetric stable laws)
\item[5.8] Theorem (domain of attraction, symmetric case)
\item[5.10] Theorem (hitting times of Brownian motion)
\item[5.11] Theorem (the waiting time of a fair gambler)
\end{results}
\subsection*{6.\ Stein's Method}
\begin{results}
\item[6.1] Lemma (Stein's lemma)
\item[6.2] Lemma (the solution and its derivatives)
\item[6.4] Theorem (CLT with an error in Wasserstein distance)
\item[6.5] Theorem (exchangeable pairs)
\item[6.7] Lemma (the Poisson Stein solution)
\item[6.8] Theorem (Chen--Stein, local dependence)
\end{results}
\subsection*{7.\ Large Deviations}
\begin{results}
\item[7.1] Lemma (the cumulant generating function)
\item[7.2] Lemma (the rate function)
\item[7.3] Theorem (Cram\'er)
\item[7.4] Lemma (the method of types)
\item[7.5] Theorem (Sanov, finite alphabet)
\item[7.6] Proposition (Gibbs' variational principle)
\item[7.7] Theorem (Gibbs conditioning principle)
\item[7.9] Theorem (Chernoff--Stein lemma)
\item[7.10] Theorem (Varadhan's lemma, for sample means)
\item[7.11] Theorem (phase transition)
\end{results}
\subsection*{8.\ Concentration of Measure}
\begin{results}
\item[8.1] Theorem (Efron--Stein)
\item[8.2] Proposition (increasing subsequences have small variance)
\item[8.3] Theorem (McDiarmid)
\item[8.4] Lemma (Ornstein--Uhlenbeck calculus)
\item[8.5] Theorem (Gaussian Poincar\'e and logarithmic Sobolev inequalities)
\item[8.6] Theorem (Gaussian concentration)
\item[8.9] Lemma (chi-square tails)
\item[8.10] Theorem (Johnson--Lindenstrauss)
\item[8.11] Lemma (convex distance as a Euclidean distance)
\item[8.12] Theorem (Talagrand's convex distance inequality)
\item[8.13] Theorem (concentration of increasing subsequences)
\end{results}
\subsection*{9.\ Martingales II: Closing the Circle}
\begin{results}
\item[9.2] Theorem (uniformly integrable martingales)
\item[9.3] Theorem (L\'evy's upward theorem, 1937)
\item[9.4] Corollary (L\'evy's zero--one law)
\item[9.5] Theorem (optional stopping for uniformly integrable martingales)
\item[9.7] Theorem (L\'evy's downward theorem)
\item[9.8] Lemma (symmetrization)
\item[9.9] Theorem (strong law of large numbers, by backward martingales)
\item[9.10] Theorem (de Finetti, 1931)
\item[9.12] Theorem (Doob, 1949)
\item[9.13] Theorem (Kakutani, 1948)
\item[9.15] Theorem (Ville's inequality, 1939)
\item[9.16] Theorem (an anytime-valid test of a coin)
\end{results}
\subsection*{10.\ Ergodic Theory}
\begin{results}
\item[10.3] Theorem (Poincar\'e recurrence, 1890)
\item[10.4] Lemma (almost invariant sets)
\item[10.5] Theorem (Kac's lemma, 1947)
\item[10.6] Theorem (von Neumann's mean ergodic theorem, 1932)
\item[10.7] Lemma (maximal ergodic lemma)
\item[10.8] Theorem (Birkhoff's ergodic theorem, 1931)
\item[10.9] Corollary (ergodicity of the examples)
\item[10.11] Proposition (Gauss)
\item[10.12] Lemma (bounded distortion)
\item[10.14] Theorem (digits of almost every number)
\item[10.15] Theorem (Roth, 1953)
\item[10.16] Theorem (Furstenberg's double recurrence)
\item[10.17] Lemma (the compact part)
\item[10.18] Lemma (almost periodicity)
\item[10.19] Lemma (the weakly mixing part)
\item[10.20] Lemma (van der Corput)
\end{results}
\subsection*{11.\ Mod-\texorpdfstring{$\varphi$}{phi} Convergence}
\begin{results}
\item[11.6] Theorem (central limit theorem)
\item[11.7] Theorem (local limit theorem with residue)
\end{results}
\subsection*{12.\ Random Permutations and the Poisson--Dirichlet Law}
\begin{results}
\item[12.1] Proposition (Cauchy's formula)
\item[12.2] Theorem (the conditioning relation)
\item[12.7] Theorem (the cycle lengths are a broken stick)
\item[12.9] Theorem (Poisson--Dirichlet limit for cycles)
\item[12.10] Theorem (Dickman's function)
\item[12.12] Corollary (the Golomb--Dickman constant)
\item[12.17] Theorem (Billingsley, 1972)
\item[12.18] Corollary (Dickman, 1930)
\end{results}
\section{Tome 11: Random Motion}
\subsection*{1.\ Gaussian Processes and the Birth of Brownian Motion}
\begin{results}
\item[1.1] Theorem (Kolmogorov's extension theorem, any index set)
\item[1.2] Corollary (Gaussian processes exist)
\item[1.5] Theorem (Kolmogorov and Chentsov)
\item[1.6] Corollary (Brownian motion, a second time)
\item[1.8] Lemma (L\'evy's maximal inequality)
\item[1.9] Lemma (maxima of Gaussians)
\item[1.10] Theorem (Paley--Wiener--Zygmund, Karhunen--Lo\`eve)
\item[1.11] Theorem (the Wiener integral)
\item[1.12] Theorem (fractional Brownian motion)
\end{results}
\subsection*{2.\ The Fine Structure of a Brownian Path}
\begin{results}
\item[2.1] Theorem (Blumenthal's zero--one law, 1957)
\item[2.4] Theorem (L\'evy, 1937)
\item[2.5] Theorem (reflection principle, joint form)
\item[2.6] Corollary (L\'evy's identity at a fixed time)
\item[2.7] Theorem (the zero set)
\item[2.9] Theorem (Taylor, 1955)
\end{results}
\subsection*{3.\ Continuous Martingales}
\begin{results}
\item[3.5] Theorem (Doob's inequalities in continuous time)
\item[3.6] Theorem (optional stopping)
\item[3.8] Theorem (convergence in $L^2$)
\item[3.14] Theorem (the quadratic variation)
\item[3.20] Corollary (hitting times)
\item[3.21] Corollary (the maximum of a drifting walk)
\item[3.22] Proposition (Wald's identities)
\end{results}
\subsection*{4.\ The It\^o Integral}
\begin{results}
\item[4.2] Lemma (Kunita--Watanabe inequality, 1967)
\item[4.6] Theorem (the It\^o integral)
\item[4.10] Proposition (Riemann sums)
\item[4.11] Theorem (It\^o's formula, 1951)
\item[4.12] Theorem (L\'evy's characterisation, 1948)
\item[4.13] Proposition (exponential martingales)
\item[4.16] Theorem (Dambis, Dubins--Schwarz, 1965)
\item[4.17] Lemma (Doob's $L^p$ inequality)
\item[4.18] Theorem (Burkholder, 1966)
\item[4.19] Theorem (Tanaka's formula)
\item[4.20] Lemma (Skorokhod's reflection lemma)
\item[4.21] Theorem (L\'evy, 1948)
\item[4.22] Corollary (the dimension of the zero set)
\end{results}
\subsection*{5.\ Brownian Motion and Analysis}
\begin{results}
\item[5.1] Theorem (the heat equation)
\item[5.2] Theorem (Feynman--Kac formula)
\item[5.4] Lemma (Bayes)
\item[5.5] Theorem (Girsanov, 1960)
\item[5.6] Corollary (Cameron--Martin formula, 1944)
\item[5.7] Lemma (the uniform law on the sphere)
\item[5.8] Proposition (strong Markov property in $\R^d$)
\item[5.9] Lemma (mean value property)
\item[5.10] Theorem (Kakutani, 1944)
\item[5.11] Theorem (Poincar\'e's cone condition)
\item[5.12] Theorem (exit from an annulus)
\item[5.14] Lemma (a common clock)
\item[5.15] Theorem (L\'evy's conformal invariance, 1948)
\item[5.16] Corollary (L\'evy)
\end{results}
\subsection*{6.\ Stochastic Differential Equations}
\begin{results}
\item[6.2] Lemma (Gronwall)
\item[6.3] Theorem (It\^o, 1951)
\item[6.5] Theorem (Markov property)
\item[6.6] Theorem (Yamada and Watanabe, 1971)
\item[6.7] Proposition (Ornstein--Uhlenbeck, 1930)
\item[6.8] Proposition (geometric Brownian motion)
\item[6.9] Theorem (Black, Scholes and Merton, 1973)
\item[6.13] Lemma (eigenvalue perturbation)
\item[6.14] Theorem (Dyson, 1962; McKean, 1969)
\item[6.15] Theorem (strong order $\frac12$)
\end{results}
\subsection*{7.\ Markov Processes and L\'evy Processes}
\begin{results}
\item[7.3] Theorem (Dynkin's martingale)
\item[7.5] Theorem (the generator of Brownian motion)
\item[7.6] Theorem (the martingale problem for $\frac12\Delta$)
\item[7.9] Proposition (exponential formula)
\item[7.12] Proposition (generator of a L\'evy process)
\item[7.13] Theorem (L\'evy--It\^o, 1934--1942)
\item[7.15] Theorem (the $\frac12$-stable subordinator)
\item[7.16] Corollary (the inverse local time)
\end{results}
\subsection*{8.\ Random Walks on Networks and Groups}
\begin{results}
\item[8.1] Lemma (maximum principle)
\item[8.2] Proposition (escape probability)
\item[8.3] Theorem (Thomson's principle)
\item[8.4] Corollary (Rayleigh's monotonicity law)
\item[8.5] Theorem (commute time identity; Chandra, Raghavan, Ruzzo, Smolensky and Tiwari, 1989)
\item[8.7] Theorem (T.~Lyons, 1983)
\item[8.8] Theorem (Nash-Williams, 1959)
\item[8.12] Lemma (Schur's test)
\item[8.13] Theorem (Kesten, 1959)
\end{results}
\subsection*{9.\ Expanders and Ramanujan Graphs}
\begin{results}
\item[9.1] Proposition (mixing)
\item[9.2] Proposition (hitting property)
\item[9.3] Theorem (expander Chernoff bound)
\item[9.4] Lemma (non-backtracking polynomials)
\item[9.5] Theorem (Kesten, 1959)
\item[9.6] Theorem (McKay, 1981)
\item[9.9] Theorem (Ihara, 1966; Bass, 1992)
\item[9.10] Theorem (Sunada, 1986)
\item[9.11] Theorem (prime number theorem for graphs)
\end{results}
\subsection*{10.\ Random Matrices I: The Semicircle}
\begin{results}
\item[10.2] Theorem (Wigner's law, almost surely)
\item[10.3] Proposition (Stieltjes inversion)
\item[10.4] Lemma (identification)
\item[10.8] Theorem (Marchenko and Pastur, 1967)
\item[10.9] Lemma (Sudakov and Fernique)
\item[10.11] Theorem (outliers of a spiked Wigner matrix)
\end{results}
\subsection*{11.\ Random Matrices II: Unitary Matrices and the Zeros of Zeta}
\begin{results}
\item[11.1] Proposition (a sum of independent Bernoulli variables)
\item[11.2] Theorem (the number variance)
\item[11.3] Corollary (Costin--Lebowitz, Soshnikov)
\item[11.4] Theorem (the sine kernel)
\item[11.5] Theorem (gap probability)
\item[11.7] Lemma (uniform vectors)
\item[11.8] Lemma (a reflection)
\item[11.9] Theorem (Bourgade--Hughes--Nikeghbali--Yor)
\item[11.10] Theorem (Keating--Snaith moments)
\item[11.11] Theorem (a central limit theorem for $\log|\det|$)
\item[11.12] Proposition (traces of powers)
\item[11.13] Theorem (Montgomery, 1973; proved in Tome 13)
\end{results}
\subsection*{12.\ Log-Correlated Fields and the Largest Values}
\begin{results}
\item[12.1] Theorem (the speed of the branching random walk)
\item[12.2] Lemma (cosine sums)
\item[12.3] Proposition (the characteristic polynomial is log-correlated)
\item[12.4] Theorem (Kahane, the $L^2$ phase)
\item[12.5] Lemma (a Bernstein inequality)
\item[12.6] Theorem (the maximum, upper bound)
\item[12.7] Theorem (Arguin--Belius--Bourgade, Paquette--Zeitouni, Chhaibi--Madaule--Najnudel; admitted)
\end{results}
\section{Tome 12: Arithmetic}
\subsection*{1.\ Number Fields and Their Integers}
\begin{results}
\item[1.16] Proposition (the sign of the discriminant)
\item[1.17] Proposition (the index formula)
\item[1.18] Theorem (Stickelberger, 1897)
\end{results}
\subsection*{2.\ Ideals and the Splitting of Primes}
\begin{results}
\item[2.4] Lemma (the missing element)
\item[2.6] Theorem (Dedekind)
\item[2.12] Theorem (Kummer, Dedekind)
\item[2.19] Theorem (splitting in cyclotomic fields)
\item[2.20] Theorem (Dedekind)
\end{results}
\subsection*{3.\ The Geometry of Numbers}
\begin{results}
\item[3.2] Lemma (Blichfeldt)
\item[3.3] Theorem (Minkowski's convex body theorem)
\item[3.7] Theorem (Minkowski's bound)
\item[3.10] Theorem (Minkowski)
\item[3.11] Theorem (Hermite)
\end{results}
\subsection*{4.\ Units}
\begin{results}
\item[4.7] Theorem (Dirichlet's unit theorem)
\item[4.10] Proposition (an Artin-type bound)
\end{results}
\subsection*{5.\ Cyclotomic Fields and Fermat}
\begin{results}
\item[5.2] Proposition (the prime $1-\zeta$)
\item[5.4] Lemma (coprime discriminants)
\item[5.6] Proposition (the Gauss sum)
\item[5.7] Theorem (quadratic reciprocity, through Frobenius)
\item[5.8] Theorem (Kummer's lemma)
\item[5.13] Theorem (Kummer, 1847)
\end{results}
\subsection*{6.\ Local Fields}
\begin{results}
\item[6.1] Theorem (Ostrowski's theorem for number fields, non-archimedean part)
\item[6.2] Proposition (the product formula)
\item[6.4] Theorem (Hensel's lemma, factorisation form)
\item[6.10] Theorem (unramified extensions)
\item[6.11] Theorem (Eisenstein)
\item[6.13] Lemma (Krasner, 1946)
\item[6.14] Proposition (roots depend continuously on the polynomial)
\end{results}
\subsection*{7.\ Quadratic Forms and the Local--Global Principle}
\begin{results}
\item[7.3] Theorem (the formulas)
\item[7.5] Theorem (Hilbert)
\item[7.7] Theorem (forms over $\Q_v$)
\item[7.8] Lemma (approximation)
\item[7.9] Theorem (existence of a rational number with given Hilbert symbols)
\item[7.10] Theorem (Hasse--Minkowski)
\item[7.12] Corollary (Meyer, 1884)
\item[7.13] Theorem (Legendre, 1785)
\item[7.14] Theorem (Gauss, Legendre)
\item[7.15] Lemma (Aubry, Davenport--Cassels)
\end{results}
\subsection*{8.\ Zeta Functions and Densities}
\begin{results}
\item[8.2] Lemma (the Lipschitz principle)
\item[8.4] Theorem (counting in a cone)
\item[8.5] Theorem (ideals in a class)
\item[8.7] Theorem (ideals in a ray class)
\item[8.9] Theorem (Dirichlet, Dedekind: the analytic class number formula)
\item[8.13] Theorem (Kummer)
\item[8.14] Lemma (the Teichm\"uller character)
\item[8.15] Lemma (Bernoulli numbers modulo $p$)
\item[8.16] Theorem (Kummer's criterion)
\item[8.20] Theorem (Kronecker, 1880)
\item[8.22] Theorem (Dirichlet's theorem over a number field)
\item[8.23] Theorem (Chebotarev, 1922)
\item[8.24] Lemma (Deuring)
\end{results}
\subsection*{9.\ Ad\`eles, Id\`eles and Tate's Thesis}
\begin{results}
\item[9.4] Proposition (Haar measure on $\Q_p$)
\item[9.6] Proposition (the Euler factors are local integrals)
\item[9.12] Theorem (adelic Poisson summation)
\item[9.18] Theorem (Tate, 1950, for $\Q$)
\item[9.19] Corollary (Riemann)
\item[9.21] Corollary (functional equation of Dirichlet $L$-functions)
\item[9.22] Lemma (adelic Minkowski)
\item[9.23] Theorem (Fujisaki)
\item[9.25] Corollary (Dirichlet's unit theorem, again)
\end{results}
\subsection*{10.\ Class Field Theory}
\begin{results}
\item[10.1] Theorem (Kronecker--Weber)
\item[10.3] Lemma (local cyclotomic fields)
\item[10.4] Lemma (tame abelian extensions)
\item[10.5] Lemma (Kummer)
\item[10.8] Proposition ($p=2$)
\item[10.9] Theorem (local Kronecker--Weber)
\item[10.10] Corollary (class field theory over $\Q$)
\item[10.11] Theorem (Artin reciprocity; admitted, not proved in this series)
\item[10.12] Theorem (existence theorem; admitted, not proved in this series)
\item[10.13] Theorem (Hilbert class field; admitted, not proved in this series)
\item[10.15] Theorem (assuming Theorem~10.13)
\item[10.19] Lemma (Gauss and Jacobi sums)
\item[10.20] Theorem (cubic reciprocity; Jacobi, Eisenstein)
\item[10.21] Theorem (Gauss)
\end{results}
\subsection*{11.\ Algebraic Curves}
\begin{results}
\item[11.4] Lemma (local decomposition)
\item[11.8] Theorem (B\'ezout)
\item[11.10] Theorem (Noether's $AF+BG$ theorem)
\item[11.15] Theorem (residue theorem)
\item[11.17] Proposition (counting forms)
\item[11.18] Proposition (Riemann's inequality)
\item[11.19] Lemma (truncated intersection numbers)
\item[11.20] Lemma (Noether's reduction lemma)
\item[11.21] Theorem (Riemann--Roch)
\item[11.22] Corollary (curves of genus 1)
\item[11.23] Proposition (Goppa codes from plane curves)
\end{results}
\subsection*{12.\ Elliptic Curves}
\begin{results}
\item[12.4] Proposition (the $p$-adic filtration)
\item[12.5] Theorem (Nagell--Lutz, 1935--1937)
\item[12.7] Theorem (Mordell, 1922)
\item[12.10] Theorem (weak Mordell--Weil)
\item[12.15] Theorem (Fermat)
\end{results}
\subsection*{13.\ Curves over Finite Fields}
\begin{results}
\item[13.4] Proposition (descent)
\item[13.6] Proposition (Euler product)
\item[13.8] Theorem (F.~K.~Schmidt; rationality and functional equation)
\item[13.9] Proposition (Gauss)
\item[13.10] Theorem (Stepanov, Bombieri)
\item[13.12] Proposition (twists)
\item[13.13] Theorem (Riemann hypothesis for $y^m=f(x)$)
\item[13.14] Corollary (Hasse, 1933)
\end{results}
\subsection*{14.\ Modular Forms and $L$-functions}
\begin{results}
\item[14.2] Theorem (valence formula)
\item[14.3] Theorem (the $j$-invariant)
\item[14.6] Theorem (Hecke)
\item[14.8] Proposition (Hecke's bound)
\item[14.9] Theorem (Hecke, 1936)
\end{results}
\section{Tome 13: Primes in Depth}
\subsection*{1.\ Zeros and Primes in Short Intervals}
\begin{results}
\item[1.1] Theorem (Vinogradov, Korobov, 1958; quoted)
\item[1.2] Theorem (Ingham, 1937)
\item[1.3] Lemma (a mean value theorem)
\item[1.5] Lemma ($\zeta$ as a short sum)
\item[1.6] Proposition (the mollified mean square)
\item[1.7] Theorem (Carlson, 1920)
\item[1.8] Corollary (Bohr and Landau, 1914)
\item[1.9] Corollary (primes in short intervals)
\end{results}
\subsection*{2.\ Siegel's Theorem}
\begin{results}
\item[2.1] Lemma (size of $L$)
\item[2.2] Lemma ($\frac{L'}L$ near the line $\sigma=1$)
\item[2.3] Theorem (Landau, Gronwall, Titchmarsh)
\item[2.4] Theorem (Landau, Page)
\item[2.5] Corollary (Page, 1935)
\item[2.8] Theorem (Siegel, 1935)
\item[2.10] Lemma (truncated Perron formula for characters)
\item[2.11] Theorem (the prime number theorem for characters)
\item[2.12] Corollary (Siegel--Walfisz, 1936)
\item[2.13] Theorem (Siegel, 1935)
\end{results}
\subsection*{3.\ Sieves and the Parity Problem}
\begin{results}
\item[3.1] Lemma (Bonferroni)
\item[3.4] Theorem (the fundamental lemma)
\item[3.5] Corollary (rough numbers)
\item[3.6] Corollary (Brun)
\item[3.7] Corollary (Brun--Titchmarsh in progressions)
\item[3.11] Theorem (Selberg's parity example)
\item[3.12] Corollary (what a sieve cannot do)
\item[3.13] Theorem (Chen, 1966 and 1973; quoted)
\item[3.17] Theorem (Lagarias--Odlyzko 1977, Serre 1981; quoted)
\item[3.18] Theorem (Hooley, 1967)
\end{results}
\subsection*{4.\ The Large Sieve and Bombieri--Vinogradov}
\begin{results}
\item[4.1] Theorem (the multiplicative large sieve)
\item[4.3] Lemma (Vaughan, 1977)
\item[4.4] Proposition (bilinear forms)
\item[4.7] Theorem (Bombieri, A.~I.~Vinogradov, 1965)
\item[4.9] Theorem (Barban; Davenport and Halberstam, 1966)
\item[4.10] Theorem (Linnik, 1941)
\end{results}
\subsection*{5.\ Exponential Sums over Primes}
\begin{results}
\item[5.1] Lemma (Dirichlet)
\item[5.3] Theorem (Vinogradov, 1937)
\item[5.4] Lemma (major arcs)
\item[5.5] Corollary (Vinogradov, 1937)
\item[5.7] Theorem (Bourgain, Demeter, Guth; Wooley for $k=3$; quoted)
\item[5.9] Lemma (moments)
\item[5.10] Theorem (Kloosterman, 1926)
\item[5.11] Theorem (Weil, 1948; quoted)
\item[5.12] Theorem (Salié, 1932)
\end{results}
\subsection*{6.\ The Circle Method I: Partitions}
\begin{results}
\item[6.2] Theorem (Hardy and Ramanujan, 1918)
\item[6.3] Theorem (Rademacher, 1937; quoted)
\end{results}
\subsection*{7.\ The Circle Method II: Three Primes and Waring's Problem}
\begin{results}
\item[7.4] Theorem (Vinogradov, 1937)
\item[7.5] Corollary (Vinogradov, 1937)
\item[7.6] Lemma (Weyl differencing)
\item[7.7] Theorem (Weyl's inequality)
\item[7.8] Theorem (Hua, 1938)
\item[7.11] Theorem (Hardy and Littlewood, 1920; Hua, 1938)
\item[7.12] Theorem (Hardy and Littlewood; quoted)
\item[7.13] Corollary (Hua, 1938)
\end{results}
\subsection*{8.\ Small Gaps between Primes}
\begin{results}
\item[8.4] Lemma (change of variables)
\item[8.7] Theorem (Maynard, 2013)
\item[8.10] Theorem (Goldston, Pintz, Yıldırım, 2005)
\item[8.11] Theorem (Zhang, 2013; quoted)
\end{results}
\subsection*{9.\ Bounded Gaps between Primes}
\begin{results}
\item[9.2] Lemma (Dirichlet)
\item[9.3] Proposition (Maynard)
\item[9.4] Theorem (Maynard, 2013)
\item[9.7] Theorem (Maynard, 2013)
\item[9.8] Theorem (Polymath, 2014; quoted)
\end{results}
\subsection*{10.\ Multiplicative Functions}
\begin{results}
\item[10.6] Theorem (Halász's inequality)
\item[10.7] Theorem (Halász, 1968)
\item[10.8] Theorem (Delange)
\item[10.10] Theorem (Wirsing, 1967)
\item[10.12] Theorem (Matomäki and Radziwiłł, 2016; quoted)
\end{results}
\subsection*{11.\ Additive Combinatorics}
\begin{results}
\item[11.2] Theorem (Freiman; Ruzsa, 1994; quoted)
\item[11.4] Lemma (Ruzsa's triangle inequality)
\item[11.5] Lemma (Ruzsa's covering lemma)
\item[11.6] Lemma (Petridis)
\item[11.7] Theorem (Plünnecke--Ruzsa inequality)
\item[11.8] Theorem (Ruzsa, 1999)
\item[11.9] Theorem (Gowers, Green, Manners and Tao, 2023; quoted)
\item[11.11] Lemma (paths of length three)
\item[11.12] Theorem (Balog--Szemerédi, 1994; Gowers, 1998)
\item[11.13] Theorem (Szemerédi and Trotter, 1983)
\item[11.14] Theorem (Elekes, 1997)
\end{results}
\subsection*{12.\ Roth's Theorem}
\begin{results}
\item[12.1] Lemma (counting lemma)
\item[12.3] Lemma (density increment)
\item[12.4] Theorem (Roth, 1953)
\item[12.5] Theorem (Behrend, 1946)
\item[12.6] Theorem (Kelley and Meka, 2023; quoted)
\item[12.7] Theorem (Szemerédi, 1975; quoted)
\end{results}
\subsection*{13.\ Gowers Norms and the Green--Tao Theorem}
\begin{results}
\item[13.3] Lemma (the $U^2$ inverse theorem)
\item[13.4] Theorem (generalized von Neumann theorem)
\item[13.7] Theorem (inverse theorem for the Gowers norms; Green, Tao and Ziegler, 2012; quoted)
\item[13.8] Theorem (Green and Tao, 2008; quoted)
\item[13.10] Theorem (relative Szemerédi theorem; Green and Tao, 2008; Conlon, Fox and Zhao, 2015; quoted)
\end{results}
\subsection*{14.\ The Critical Line}
\begin{results}
\item[14.2] Lemma (theta near $i$)
\item[14.3] Theorem (Hardy, 1914)
\item[14.4] Theorem (Selberg, 1942; Levinson, 1974; Conrey, 1989; Pratt, Robles, Zaharescu and Zeindler, 2020; quoted)
\item[14.5] Lemma (first derivative test for integrals)
\item[14.6] Lemma ($\zeta$ as a sum of length $t$)
\item[14.7] Theorem (Hilbert's inequality, with Schur's constant)
\item[14.9] Theorem (Hardy and Littlewood, 1918)
\item[14.10] Theorem (Hardy and Littlewood, 1921; quoted)
\item[14.11] Theorem (Riemann, Siegel, 1932; quoted)
\item[14.13] Lemma (Montgomery's explicit formula)
\item[14.15] Theorem (Montgomery, 1973)
\item[14.16] Corollary (Montgomery, 1973)
\end{results}
\section{Tome 14: Random Arithmetic}
\subsection*{1.\ Additive Functions and the Kubilius Model}
\begin{results}
\item[1.2] Theorem (Tur\'an--Kubilius)
\item[1.3] Theorem (Kubilius)
\item[1.5] Theorem (Erd\H{o}s and Wintner, 1939)
\item[1.6] Lemma (Littlewood--Offord, in Erd\H{o}s's form, 1945)
\item[1.7] Theorem (L\'evy, 1931)
\item[1.9] Corollary (Schoenberg 1928, Davenport 1933)
\item[1.10] Theorem (Jessen and Wintner, 1935)
\end{results}
\subsection*{2.\ The Selberg--Delange Method}
\begin{results}
\item[2.1] Lemma (Hankel, 1864)
\item[2.5] Theorem (Selberg--Delange)
\item[2.7] Corollary (Selberg, 1954)
\item[2.9] Corollary (Landau, 1908)
\end{results}
\subsection*{3.\ The Erd\H{o}s--Kac Theorem, Sharply}
\begin{results}
\item[3.2] Theorem (Sathe, 1953; Selberg, 1954)
\item[3.3] Theorem (Rényi and Turán, 1958)
\end{results}
\subsection*{4.\ The Anatomy of Integers}
\begin{results}
\item[4.1] Proposition (Rankin's trick)
\item[4.5] Lemma (Dixon)
\item[4.6] Theorem (Dixon, 1981)
\item[4.8] Theorem (Deshouillers, Dress and Tenenbaum, 1979)
\item[4.9] Theorem (Erd\H{o}s, 1960)
\end{results}
\subsection*{5.\ Random Multiplicative Functions}
\begin{results}
\item[5.4] Theorem (Wintner, 1944)
\end{results}
\subsection*{6.\ The Bohr--Jessen Law}
\begin{results}
\item[6.2] Lemma (Kronecker--Weyl for the primes)
\item[6.5] Theorem (Bohr and Jessen, 1930--1932)
\item[6.7] Corollary (Bohr and Courant, 1914)
\item[6.8] Theorem (Jessen and Wintner, 1935)
\end{results}
\subsection*{7.\ Selberg's Central Limit Theorem}
\begin{results}
\item[7.1] Theorem (Selberg, 1946)
\item[7.14] Theorem (a twisted second moment)
\end{results}
\subsection*{8.\ Moments of Zeta and Random Matrices}
\begin{results}
\item[8.1] Theorem (Hardy and Littlewood, 1918; Ingham, 1926)
\item[8.3] Theorem (Ramachandra; Heath-Brown; Rudnick and Soundararajan)
\end{results}
\subsection*{9.\ Chebyshev's Bias}
\begin{results}
\item[9.5] Lemma (Kronecker--Weyl for real frequencies)
\item[9.6] Theorem (Wintner, 1935; Rubinstein and Sarnak, 1994)
\item[9.7] Theorem (Rubinstein and Sarnak, 1994)
\end{results}
\subsection*{10.\ Primes to Large Moduli}
\begin{results}
\item[10.2] Proposition (the smoothed explicit formula)
\item[10.3] Lemma (the sieve bound)
\item[10.5] Lemma (duality)
\item[10.6] Lemma (smooth character sums)
\item[10.7] Lemma (a zero detector)
\item[10.8] Theorem (log-free zero density; Linnik, 1944; Gallagher, 1970)
\item[10.9] Lemma (one exceptional zero)
\item[10.10] Theorem (Gallagher, 1970)
\item[10.11] Corollary (Linnik's theorem without exceptional zero)
\end{results}
\subsection*{11.\ Gaps and Clumps}
\begin{results}
\item[11.3] Theorem (Buchstab, 1937; de~Bruijn, 1950)
\item[11.6] Theorem (Maier, 1985)
\item[11.7] Lemma (good moduli)
\item[11.10] Theorem (Erd\H{o}s, 1935; Rankin, 1938)
\end{results}
\subsection*{12.\ Equidistribution and Random Polynomials}
\begin{results}
\item[12.1] Theorem (an Erd\H{o}s--Turán inequality)
\item[12.3] Proposition (Birch, 1968)
\item[12.5] Lemma (Jordan, 1872)
\item[12.7] Theorem (van der Waerden, 1936)
\end{results}
\subsection*{13.\ Random Groups and the Cohen--Lenstra Heuristics}
\begin{results}
\item[13.5] Theorem (moments of random cokernels)
\item[13.6] Theorem (Friedman--Washington, 1989)
\item[13.7] Corollary (the Cohen--Lenstra law)
\item[13.8] Lemma (Wood, 2017)
\item[13.9] Theorem (Wood, 2017)
\end{results}
\section{Tome 15: The Research Workshop}
\subsection*{1.\ How Research Works}
\begin{results}
\item[1.2] Proposition (the counting bound)
\item[1.3] Proposition (doubling)
\item[1.4] Theorem (Erd\H os and Moser, in a sharpened form)
\item[1.6] Proposition (from $Q$-dissociated to dissociated)
\item[1.7] Proposition (lattices give $Q$-dissociated sequences)
\item[1.8] Lemma (the odd cycle)
\end{results}
\subsection*{2.\ How to Read a Paper}
\begin{results}
\item[2.1] Lemma (Moser)
\item[2.2] Proposition (Berry--Esseen for signed sums)
\item[2.3] Theorem (Dubroff, Fox and Xu; first proof)
\item[2.4] Proposition (the reduction to the boundary)
\end{results}
\subsection*{3.\ Reading a Classic: Shuffling Cards}
\begin{results}
\item[3.1] Theorem (Bayer and Diaconis, 1992)
\item[3.2] Lemma (Tanny, 1973)
\item[3.5] Corollary (cutoff at $\frac32\log_2n$)
\end{results}
\subsection*{4.\ Guided Reading: Selberg's Theorem at Many Points}
\begin{results}
\item[4.1] Lemma (the covariance of the prime sums)
\item[4.2] Theorem (a multi-point central limit theorem)
\item[4.3] Lemma (joint moments)
\end{results}
\subsection*{5.\ Guided Reading: Random Groups}
\begin{results}
\item[5.1] Lemma (counting surjections)
\item[5.2] Theorem (moments of symmetric cokernels)
\item[5.3] Lemma (one more row and column)
\item[5.4] Theorem (the corank law)
\end{results}
\subsection*{6.\ Guided Reading: Better Than Square Root}
\begin{results}
\item[6.2] Lemma (many-to-one)
\item[6.3] Lemma (ballot bound)
\item[6.4] Lemma (no high individuals)
\item[6.5] Theorem (better than square root, in the tree)
\end{results}
\subsection*{7.\ From Data to Conjecture}
\begin{results}
\item[7.1] Theorem (Bailey, Borwein and Plouffe, 1995)
\item[7.2] Proposition (digit extraction)
\item[7.5] Theorem (Lenstra, Lenstra and Lov\'asz, 1982)
\item[7.6] Proposition (what the search guarantees)
\end{results}
\subsection*{8.\ Heuristics That Predict}
\begin{results}
\item[8.2] Proposition (Verebrusov, 1908; Mahler, 1936)
\item[8.4] Proposition (the divisor trick)
\end{results}
\subsection*{9.\ Proofs a Computer Can Check}
\begin{results}
\item[9.3] Lemma (uniqueness of types)
\item[9.4] Theorem (Curry--Howard, implicational fragment)
\item[9.7] Theorem (soundness)
\end{results}
\subsection*{10.\ Writing It Down}
\begin{results}
\item[10.1] Theorem (Jordan's theorem for polygons)
\end{results}
\subsection*{11.\ Problems to Work On}
\begin{results}
\item[11.1] Proposition (Singmaster, 1971)
\item[11.2] Proposition (Singmaster, 1975)
\item[11.5] Proposition (Rosa, 1967)
\item[11.6] Lemma (Kelly's lemma)
\item[11.8] Proposition (Lehmer, 1932)
\item[11.10] Proposition (Gauss--Lucas)
\end{results}
\subsection*{12.\ Five Projects}
No named results; 4 exercises.
